\section{Minimum-Benefit Certification: Complete Derivations}
\label{app:minbenefit}

This appendix separates two routes. The general-objective route uses a
numerically bounded asymptotic remainder. The second route uses the
variational structure of three-dimensional, zero-Dirichlet Poisson compliance
to bound a finite insertion directly. The latter needs no unknown remainder
constant and permits Lipschitz outer geometry. It does not certify the
previously reported ranking runs retrospectively; \cref{app:applicationcertificate} evaluates a separate guarded implementation. All certificates below
concern the stated continuum PDE under their oracle assumptions, not a
fluid-cooling model or unbounded floating-point geometry error.

\subsection{What is being guaranteed?}

Write the actual benefit of an insertion as
\[
 B(z,r)=J(D)-J(D\setminus\overline B_r(z)).
\]
A lower bound $\underline B(z,r)$ is useful when it exceeds a specified
target $\tau\ge0$. The statistical statement is
\begin{equation}
 \Prob\{\text{accept and }B(z,r)<\tau\}\le\delta,
 \label{eq:benefitfailure}
\end{equation}
including any covered data-dependent choice of radius or target. It is not
$\Prob(B\ge\tau\mid\text{accept})\ge1-\delta$: conditioning on a rare
acceptance can change that probability. In ideal Dirichlet cooling,
\cref{prop:coolingmonotonicity} already supplies nonnegative benefit.
Minimum-benefit certification asks the additional quantitative question.

\begin{proof}[Proof of \cref{thm:minbenefit}]
\textbf{Step 1: one event bounds the coefficient.}
Let $E$ be the event in \cref{eq:validationcoverage}. Set
$b=b_{\mathrm{term}}+b_{\mathrm{cache}}$ and
$s=\widehat{\DT}_N+r_N(\delta)+b$. The bias assumption gives
$\DT\le\E\widehat{\DT}_N+b$, and on $E$ the confidence bound gives
$\E\widehat{\DT}_N\le\widehat{\DT}_N+r_N$. Thus $\DT\le s$ on $E$.

\textbf{Step 2: translate derivative uncertainty into objective units.}
The upper side of the deterministic remainder inequality gives
\[
 J(D_r)-J(D)\le\rho(r)\DT+C_z\rho(r)^{1+\alpha}.
\]
Negating both sides and substituting $\DT\le s$ proves, on $E$,
\[
 B(z,r)\ge-\rho(r)s-C_z\rho(r)^{1+\alpha}
             =\underline B_{\rm TD}(z,r).
\]
Notice the factor $\rho$: a confidence bound on a derivative alone does not
have the units of an actual objective improvement.

\textbf{Step 3: radii and acceptance.}
For fixed $z$, the same $E$ controls every radius for which the deterministic
remainder bound holds uniformly. Evaluating the displayed inequality at a
measurable selected radius preserves it. On $E$, accepting only when
$\underline B_{\rm TD}\ge\tau$ implies $B\ge\tau$, even for a measurable
data-dependent $\tau$. Therefore every false acceptance lies in $E^c$, and
its probability is at most $\delta$. This reasoning allows rejection on all
outcomes and does not presume that a feasible certified radius exists.
\end{proof}

\paragraph{Capacity must also be bounded.}
The preceding expression uses the actual relative operator capacity, not
merely its leading asymptotic equivalent. If only certified numerical bounds
$0<\rho_-\le\rho\le\rho_+$ are available, a conservative replacement is
\begin{equation}
 \underline B_{\rm int}=
 \begin{cases}
 -\rho_-s-C_z\rho_+^{1+\alpha},&s<0,\\
 -\rho_+s-C_z\rho_+^{1+\alpha},&s\ge0.
 \end{cases}
 \label{eq:capacityintervalbenefit}
\end{equation}
Indeed, $-\rho s$ is increasing in $\rho$ when $s<0$ and decreasing when
$s\ge0$, whereas $-C_z\rho^{1+\alpha}\ge-C_z\rho_+^{1+\alpha}$ in both
cases. Adding these two lower bounds proves the claim. Replacing $\rho$ by
$4\pi r$ without an error allowance is generally invalid in this route.

\subsection{A direct finite-hole bound for thermal compliance}

Throughout the rest of this appendix, $D\subset\R^3$ is bounded and
Lipschitz, $f\in L^\infty(D)$ satisfies $0\le f\le F$, and
$u\in H_0^1(D)$ solves $-\Delta u=f$. For $\overline B_R(z)\Subset D$,
let known constants $f_-,f_+$ satisfy
$0\le f_-\le f(x)\le f_+$ almost everywhere on $B_R(z)$.
Set $K=\overline B_r(z)$ with $0<r<R$, and let $v$ solve the same PDE on
$D\setminus K$ with zero trace on every boundary component. Extend $v$ by
zero on $K$. The objective is the unnormalized compliance $J(D)=\int_Dfu$.
Constant diffusivity has been scaled to one. Different diffusivities require
the corresponding factors in both the PDE and energy.

\begin{lemma}[Positivity, boundedness, and an interior representative]
\label{lem:benefitstatebounds}
The state $u$ is nonnegative. If $\overline D\subset B_L(c)$, then
\begin{equation}
 0\le u(x)\le F(L^2-|x-c|^2)/6\le FL^2/6
 \label{eq:enclosingballbound}
\end{equation}
The inequalities hold almost everywhere in $D$. The state has a continuous representative on every compact subset of $D$.
\end{lemma}
\begin{proof}
\textbf{Step 1: existence and the sign.}
Poincare and Cauchy--Schwarz make $\int_D\nabla u\cdot\nabla\phi=\int_Df\phi$
a coercive weak problem; Lax--Milgram gives its unique solution.
The negative part $u^-=\max(-u,0)$ belongs to $H_0^1(D)$. Testing with it
yields $-\int_D|\nabla u^-|^2=\int_Dfu^-\ge0$. Both sides must vanish;
Poincare implies $u^-=0$.

\textbf{Step 2: an explicit supersolution.}
Put $Q(x)=F(L^2-|x-c|^2)/6$. Since $\Delta|x-c|^2=6$ in three dimensions,
$-\Delta Q=F\ge f$. The trace of $u-Q$ is nonpositive. Consequently
$(u-Q)^+\in H_0^1(D)$ is an admissible test in the difference of the two
weak equations and gives
\[
 \int_D|\nabla(u-Q)^+|^2=\int_D(f-F)(u-Q)^+\le0.
\]
Poincare implies $u\le Q$. This proves \cref{eq:enclosingballbound}.

\textbf{Step 3: point values.}
Interior elliptic regularity for the Laplacian and bounded $f$ gives
$u\in W_{\rm loc}^{2,p}(D)$ for finite $p$; take $p>3$ and apply the local
Sobolev embedding to obtain a continuous (in fact locally $C^{1,\gamma}$
for some $\gamma>0$) representative. These standard local regularity and
trace results are the same background facts used in
\cref{app:deterministic}~\cite{evans2010pde,gilbarg2001elliptic}.
The almost-everywhere bounds hold everywhere for this representative in
the interior.
\end{proof}

\begin{lemma}[From a cavity temperature floor to a finite benefit]
\label{lem:floorcapacity}
If $u\ge m\ge0$ on $K$, then
\begin{equation}
 J(D)-J(D\setminus K)\ge4\pi r m^2.
 \label{eq:floorcapacity}
\end{equation}
\end{lemma}
\begin{proof}
\textbf{Step 1: the exact finite-domain identity.}
The zero extension of $v$ belongs to $H_0^1(D)$ and its gradient vanishes
almost everywhere on $K$. Testing the state equations gives
$\int_D|\nabla u|^2=J(D)$,
$\int_D|\nabla v|^2=J(D\setminus K)$, and
$\int_D\nabla u\cdot\nabla v=J(D\setminus K)$.
Expanding the square therefore gives
\[
 J(D)-J(D\setminus K)=\int_D|\nabla w|^2,\qquad w=u-v.
\]
This is \cref{eq:coolingprojection}; it is exact for finite $r$.

\textbf{Step 2: turn $w$ into a capacity competitor.}
If $m=0$, the energy identity already proves the result. Otherwise extend
$w/m$ by zero from $D$ to $\R^3$. It belongs to the homogeneous energy space
$\mathcal D^{1,2}(\R^3)$, defined as the completion of
$C_c^\infty(\R^3)$ in gradient norm. It is at least one on $K$ in the
quasi-everywhere Sobolev sense: $w=u$ there, the zero extension of $v$
vanishes quasi-everywhere there, and the continuous $u$ is at least $m$.
By the definition of Newtonian capacity,
\[
 m^{-2}\int_D|\nabla w|^2\ge\operatorname{cap}_{\R^3}(K).
\]

\textbf{Step 3: fix the capacity normalization.}
The equilibrium potential is $P(x)=\min(1,r/|x-z|)$. Its exterior energy is
$4\pi\int_r^\infty r^2s^{-2}\,ds=4\pi r$, and its interior energy is
zero. To verify that no smaller energy is possible, truncate any admissible
competitor to the interval $[0,1]$; this cannot increase gradient energy and
sets it to one on the ball. For a smooth compactly supported such competitor
$a$, its spherical average $A(s)$ satisfies $A(r)=1$, $A(\infty)=0$.
Jensen and Cauchy--Schwarz give
\begin{align*}
 \int_{|x-z|>r}|\nabla a|^2
 &\ge4\pi\int_r^\infty s^2|A'(s)|^2\,ds,\\
 1=\left|\int_r^\infty A'(s)\,ds\right|^2
 &\le\left(\int_r^\infty s^2|A'(s)|^2\,ds\right)
       \left(\int_r^\infty s^{-2}\,ds\right).
\end{align*}
The last factor is $1/r$, so the energy is at least $4\pi r$.
Approximation in the exterior energy space with fixed trace extends this
inequality to the admissible truncated Sobolev competitors. Equivalently,
subtract $P$ and use the zero-trace energy-density argument of
\cref{prop:budgetgeometry}. Thus the infimum equals $4\pi r$.
Substituting into Step 2 proves the claimed finite-hole bound.
\end{proof}

\begin{lemma}[Explicit temperature floor from one interior value]
\label{lem:balltemperaturefloor}
Define $t=r/R$ and
\begin{equation}
 \begin{split}
 c(t)&=(1-t)/(1+t)^2,\\
 m(L_u;r,R)&=\frac{f_-(R^2-r^2)}6\\
 &\quad+c(t)\left(L_u-\frac{f_+R^2}6\right)_+,
 \end{split}
 \label{eq:temperaturefloor}
\end{equation}
where $(a)_+=\max(a,0)$. For $L_u\le u(z)$, this gives the floor
$u(x)\ge m(L_u;r,R)$ throughout $\overline B_r(z)$.
\end{lemma}
\begin{proof}
\textbf{Step 1: separate the source and harmonic parts inside the ball.}
Let $q\in H_0^1(B_R(z))$ solve $-\Delta q=f$ there. Define
$h=u-q$ on that ball. Subtracting the weak equations gives $\Delta h=0$.
Its trace is $u\ge0$, so the negative-part test proves $h\ge0$.
The polynomial functions
$q_\pm(x)=f_\pm(R^2-|x-z|^2)/6$ vanish on the sphere and solve
$-\Delta q_\pm=f_\pm$. The positive-part comparison argument in
\cref{lem:benefitstatebounds} applied to $q-q_+$ and $q_--q$ gives
$q_-\le q\le q_+$. In particular,
\[
 h(z)=u(z)-q(z)\ge u(z)-f_+R^2/6,
 \qquad h(z)\ge0.
\]
Hence $h(z)\ge(L_u-f_+R^2/6)_+$.

\textbf{Step 2: derive the harmonic comparison constant.}
Translate $z$ to zero. For $|x|=s<R$, the Poisson kernel representation is
\[
 h(x)=\int_{|y|=R}\frac{R^2-s^2}{4\pi R|x-y|^3}h(y)\,dS_y.
\]
At the center the kernel is $1/(4\pi R^2)$, so its pointwise ratio to the
central kernel obeys
\[
 \frac{R(R^2-s^2)}{|x-y|^3}
 \ge\frac{R(R^2-s^2)}{(R+s)^3}
 =\frac{1-s/R}{(1+s/R)^2}.
\]
Multiplication by $h(y)\ge0$ and integration therefore imply
$h(x)\ge c(s/R)h(0)$. Finally
$c'(a)=(a-3)/(1+a)^3<0$ for $0\le a<1$; thus $s\le r$ gives
$h(x)\ge c(r/R)h(z)$. The Poisson formula is valid with the continuous
boundary values here because $\overline B_R(z)$ lies strictly inside $D$;
the zero-Dirichlet bounded-source solution $q$ is continuous up to that sphere.

\textbf{Step 3: add the two lower bounds.}
On $\overline B_r(z)$,
$q(x)\ge f_-(R^2-r^2)/6$. Adding the lower bound on $h$ from Steps 1--2
gives exactly \cref{eq:temperaturefloor}. Both summands are nonnegative.
This proof uses a finite interior ball, not a small-$r$ Taylor approximation.
\end{proof}

\begin{proof}[Proof of \cref{thm:coolingbenefit}]
Work on the event $E_u=\{L_u\le u(z)\}$. Then \cref{lem:balltemperaturefloor} supplies a
nonnegative temperature floor on the entire inserted ball. Substitution into
\cref{lem:floorcapacity} proves
$B(z,r)\ge4\pi r\,m(L_u;r,R)^2=\underline B_{\rm cool}(z,r)$.
Both lemmas hold for every admissible $0<r<R$ simultaneously on the same
$E_u$. They also hold for every admissible $R$ if its source bounds and
geometric containment are valid. Thus choosing a radius after observing
$L_u$ does not require a union bound over radii. The event
$\{\text{accept},B<\tau\}$ is contained in $E_u^c$, proving its probability
bound. In particular, $L_u=0$ is deterministically valid and yields
\begin{equation}
 \underline B_{\rm source}(z,r)
 =\frac{\pi r f_-^2(R^2-r^2)^2}{9}.
 \label{eq:sourceonlybenefit}
\end{equation}
This source-only bound may be small or zero; its validity does not assert
that it can meet the requested target. A better certified $L_u$ can improve
the bound because $m$ is nondecreasing in $L_u$.
\end{proof}

\paragraph{Relation to the topological score.}
For an exact value $L_u=u(z)$ and constant $f=f_0$ on $B_R(z)$,
$u(z)\ge f_0R^2/6$ by comparison. Sending $r\downarrow0$ in
\cref{eq:temperaturefloor} gives $m\to u(z)$ and hence
$\underline B_{\rm cool}/(4\pi r)\to u(z)^2=-D_T(z)$.
For merely bounded nonconstant $f$, one can instead choose interior radii
$R(r)\to0$ with $r/R(r)\to0$ and uniformly bounded $f_+$; then $c\to1$
and both source terms vanish, giving the same limit. These statements follow
by direct substitution and continuity of $(\cdot)_+$, using $u(z)\ge0$.
A conservative or noisy $L_u$ need not attain this limiting coefficient.
The direct finite-radius bound builds on classical energy, capacity, and
harmonic comparison; we do not claim these constituent results as new.

\subsection{A one-sided WoS validator with an explicit range bound}

An unbiased two-sided state estimator is not necessary for a lower
temperature bound. Nonnegative source contributions permit a deliberately
truncated estimator whose omitted terminal value has a known sign. This
section specifies a separate validator, not an assumption about the existing
application ranking code.

\begin{proposition}[Finite-walk lower confidence bound]
\label{prop:onesidedwos}
Fix integers $T\ge1$, $N\ge1$, a radius cap $s_*>0$, and
$\delta\in(0,1)$ before observing the validation data. Start $X_0=z$.
At each active step $k<T$, choose a history-measurable
$0<s_k\le s_*$ with $\overline B_{s_k}(X_k)\Subset D$. Draw a uniform
spherical direction $\Theta_k$ and set $X_{k+1}=X_k+s_k\Theta_k$.
Also draw $Y_k$ in that ball with density
\[
 \frac{6}{s_k^2}G_{B_{s_k}(X_k)}(X_k,y),\qquad
 S_k=\frac{s_k^2}{6}f(Y_k).
\]
All fresh draws have their stated conditional distributions given the past;
the source and exit draws may be independent. Stop at any history-measurable
rule by step $T$. After stopping, set $S_k=0$ and keep $X_k$ fixed. Define
$W=\sum_{k=0}^{T-1}S_k$. For independent identically distributed complete
walks $W_1,\ldots,W_N$, set
\begin{equation}
 A=\frac{TFs_*^2}{6},\qquad
 L_u=\max\left(0,\overline W-
 A\sqrt{\frac{\log(1/\delta)}{2N}}\right).
 \label{eq:onesidedwos}
\end{equation}
Then $\Prob\{L_u\le u(z)\}\ge1-\delta$. For $F=0$, the solution and
estimator are identically zero and the formula is read with $A=0$.
\end{proposition}
\begin{proof}
\textbf{Step 1: the one-ball identity and normalization.}
On a strictly contained ball, decompose $u$ into the zero-boundary source
solution and the harmonic extension of its boundary values. At the center,
the Poisson kernel is uniform on the sphere, so
\[
 u(X_k)=\E[u(X_{k+1})\mid\mathcal G_k]
      +\int_{B_{s_k}(X_k)}G(X_k,y)f(y)\,dy,
\]
where $\mathcal G_k$ is the history before the step. The kernel is
$G(X_k,y)=(4\pi)^{-1}(|y-X_k|^{-1}-s_k^{-1})$. Spherical integration gives
\[
 \int G(X_k,y)\,dy
 =\int_0^{s_k}(a-a^2/s_k)\,da=s_k^2/6.
\]
Consequently the stated density integrates to one and
$\E[S_k\mid\mathcal G_k]=\int Gf$.
An implementable sampler is
$Y_k=X_k+s_k Q_k\Phi_k$, where $\Phi_k$ is uniform on $\mathbb S^2$ and
$Q_k$ has density $6q(1-q)$ on $(0,1)$, as derived in
\cref{prop:greenballsample}. For constant $f$, $S_k=f s_k^2/6$ is exact
without drawing $Y_k$.

\textbf{Step 2: telescope at a bounded stopping time.}
At an active step the preceding identity is
$\E[S_k+u(X_{k+1})\mid\mathcal G_k]=u(X_k)$.
At a stopped step it remains true because $S_k=0$ and $X_{k+1}=X_k$.
Each term is integrable: \cref{lem:benefitstatebounds} bounds $u$ and
$0\le S_k\le Fs_*^2/6$. Taking expectations and summing all $T$ finite
identities cancels consecutive state values and gives
\begin{equation}
 u(z)=\E W+\E u(X_T),\qquad 0\le\E W\le u(z).
 \label{eq:truncationidentity}
\end{equation}
Thus shell termination and a finite path cap create downward bias, whose
sign is sufficient for this one-sided use. No assertion of zero bias is made.

\textbf{Step 3: prove the needed bounded-variable concentration.}
For any $Z\in[0,A]$, let
$g(\lambda)=\log\E\exp(\lambda(Z-\E Z))$.
Boundedness permits differentiating under the expectation. Under the
exponentially tilted distribution, $g''(\lambda)$ is the variance of $Z$.
Every distribution supported on $[0,A]$ has variance at most $A^2/4$:
$\Var Z=\min_c\E(Z-c)^2\le\E(Z-A/2)^2\le A^2/4$.
Since $g(0)=g'(0)=0$, integrating the second-derivative bound twice gives
$g(\lambda)\le\lambda^2A^2/8$ for $\lambda\ge0$.
Independence of the $N$ walks, Markov's inequality, and this bound yield
\[
 \Prob\{\overline W-\E W\ge a\}
 \le\exp(-\lambda Na+N\lambda^2A^2/8).
\]
For $A>0$ and $a>0$, the choice $\lambda=4a/A^2$ gives the bound
$\exp(-2Na^2/A^2)$. With $a$ equal to the sampling allowance in
\cref{eq:onesidedwos}, this is $\delta$.
This supplies the one-sided Hoeffding inequality used here with all constants.

\textbf{Step 4: convert it into a state bound.}
On the complementary event, $\overline W-a\le\E W\le u(z)$ by Step 2.
Since $u(z)\ge0$, replacing $\overline W-a$ by its maximum with zero
preserves the inequality. This is precisely \cref{eq:onesidedwos}.
\end{proof}

\paragraph{What bias has and has not been handled.}
In \cref{eq:truncationidentity}, the exact truncation bias is
$u(z)-\E W=\E u(X_T)\in[0,FL^2/6]$. Its sign makes a separate magnitude
bound unnecessary for the proposed lower certificate. This does not cover
an approximate source integral that overestimates $f$, an unsafe distance
query that steps outside $D$, approximate boundary data, or a correlated
collection treated as independent walks. A rigorously known lower source
$\widetilde f\in[0,f]$ may be substituted: the one-ball inequality becomes
$\E[S_k+u(X_{k+1})\mid\mathcal G_k]\le u(X_k)$ and telescoping still gives
$\E W\le u(z)$. Alternatively, if a numerical walk output $\widetilde W$
has a certified overshoot bound $\E\widetilde W\le u(z)+\beta$ and lies in
a known interval of width $A$, subtracting $\beta$ as well as the sampling
radius proves the same result. An unquantified numerical overshoot is not
covered. Increasing $T$ need not tighten the certificate because the range
bound $A$ grows with $T$; this is a conservative construction, not an
efficiency claim.

\subsection{Costs, relative targets, selection, and repeated insertions}

\begin{proposition}[Explicit costs and fixed normalization]
\label{prop:benefitcost}
Suppose $B\ge\underline B$ on a coverage event. If the added design cost
obeys $\Delta C\le c_+$, the net improvement in $J+C$ obeys
\[
 (J+C)_{\rm old}-(J+C)_{\rm new}\ge\underline B-c_+.
\]
Thus $\underline B\ge\tau+c_+$ certifies net benefit at least $\tau$ on
that event. For the objective $J/V_{\rm ref}+C$ with fixed known
$V_{\rm ref}>0$, the corresponding bound is $\underline B/V_{\rm ref}-c_+$.
For a relative target $B/J_0\ge\eta\ge0$, suppose $J_0>0$ and a certified
upper bound $J_0\le U_0$ is available on a second event. Then
$\underline B\ge\eta U_0$ suffices on the intersection of both events.
\end{proposition}
\begin{proof}
\textbf{Step 1: subtract costs in the correct units.}
By definition, the net objective reduction is $B-\Delta C$.
The inequalities $B\ge\underline B$ and $-\Delta C\ge-c_+$ give the first
claim by addition. Dividing only the compliance term by fixed
$V_{\rm ref}>0$ preserves its order and gives the normalized claim.

\textbf{Step 2: justify an explicit geometric cost example.}
For one disjoint interior sphere, a stipulated cost
$C=\lambda V_{\rm removed}+\gamma A_{\rm cavity}+\kappa n$ with nonnegative
coefficients increases by
\[
 c_+=\lambda\frac{4\pi r^3}{3}+\gamma4\pi r^2+\kappa.
\]
The ball volume, newly exposed sphere area, and one-unit increase of count
give these terms exactly. The formula is not asserted for overlapping
balls or for an unspecified manufacturing cost model.

\textbf{Step 3: a relative threshold uses an upper denominator bound.}
On the intersection of the two events,
$B\ge\underline B\ge\eta U_0\ge\eta J_0$.
Division by $J_0>0$ proves the relative claim. If the events fail with
probabilities at most $\delta$ and $\delta_0$, their intersection fails
with probability at most $\delta+\delta_0$ by subadditivity; independence
is not necessary. Plugging in an unbounded noisy estimate of $J_0$ does not
supply this implication.
\end{proof}

\begin{proposition}[Adaptive minimum-benefit guarantees]
\label{prop:adaptivebenefit}
At round $t$, condition on the complete past $\mathcal F_{t-1}$, including
all earlier geometry changes. Choose a finite shortlist of $m_t$ candidates
before collecting their validation data. Allocate nonnegative,
$\mathcal F_{t-1}$-measurable error levels $\delta_{tj}$ with
$\sum_{t,j}\delta_{tj}\le\delta_{\rm total}$ almost surely.
Assume every candidate certificate has conditional coverage at least
$1-\delta_{tj}$ and all its deterministic hypotheses hold on the current
domain. Select at most one insertion in the round, accepting it only if
its covered lower benefit (or net benefit) is at least its target $\tau_t$.
Then the probability that any accepted insertion misses its target is at
most $\delta_{\rm total}$. On the same simultaneous event, a finite sequence
of accepted insertions guarantees total benefit at least the sum of their
targets, provided no other unaccounted geometry changes occur.
\end{proposition}
\begin{proof}
\textbf{Step 1: cover the shortlist before choosing its winner.}
Let $E_{tj}$ be candidate $j$'s conditional coverage event. On
$E_t=\bigcap_{j=1}^{m_t}E_{tj}$, every displayed candidate bound is true
at once. Any winner selected from these data is therefore covered, as is
any radius selected using its uniform radius guarantee. Conditional
subadditivity gives
$\Prob(E_t^c\mid\mathcal F_{t-1})\le\sum_j\delta_{tj}$.
Candidates may share randomness if their marginal conditional bounds remain
valid. The data used to invent the shortlist cannot silently be reused as
fresh validation data.

\textbf{Step 2: control all rounds.}
A false acceptance at round $t$ is a subset of $E_t^c$. Taking expectations,
then applying countable subadditivity and Tonelli to nonnegative terms gives
\[
 \Prob(\text{some false acceptance})
 \le\sum_t\E\sum_j\delta_{tj}
 =\E\sum_{t,j}\delta_{tj}\le\delta_{\rm total}.
\]
Recomputing the PDE and containment data on the updated domain is essential:
each certificate refers to that domain, not the original unperforated object.

\textbf{Step 3: telescoping actual finite benefits.}
For accepted designs $D_0,D_1,\ldots,D_k$ and their actual objective values,
$J(D_{i-1})-J(D_i)\ge\tau_i$ holds simultaneously on the covered event.
Summing cancels the intermediate objectives and yields
$J(D_0)-J(D_k)\ge\sum_{i=1}^k\tau_i$. The same algebra applies to $J+C$.
Unvalidated shape refinement between those designs would contribute an
additional objective change and cannot be omitted from this sum.

\textbf{Step 4: sample-size and stopping choices.}
The preceding proof also covers stages indexed by $(t,j,\ell)$ if a
predeclared sequence of fixed sample sizes is inspected with error levels
whose total obeys the same budget. Each stage needs its own valid marginal
coverage; the batches may be nested because subadditivity does not require
independence across stages. A single fixed-$N$ bound repeatedly checked
until it passes supplies no such guarantee. Within each use of
\cref{prop:onesidedwos}, independence is still required across its complete
walks. Selecting one site using scouting data and then drawing a single
fresh final validation batch is the special case $m_t=1$.
\end{proof}

\subsection{An exactly solvable finite-radius check}
\label{app:benefitcheck}

\begin{proposition}[Concentric unit-ball reference]
\label{prop:benefitballreference}
For $D=B_1(0)$, $f=1$, and a concentric radius-$r$ cavity with $0<r<1$,
write $s=|x|$. The original and perforated states are
\begin{align}
 u(s)&=(1-s^2)/6,\\
 v(s)&=(1-s^2)/6-\frac{r(1+r)}6(1/s-1),\quad r<s<1.
\end{align}
Their exact compliance benefit is
\begin{equation}
 B_{\rm exact}(r)=\frac{\pi r(1-r)(1+r)^2}{9}
                       +\frac{4\pi r^5}{45},\qquad J_0=\frac{4\pi}{45}.
 \label{eq:exactballbenefit}
\end{equation}
The direct capacity bound using the exact temperature floor is
$B_{\rm floor}=\pi r(1-r^2)^2/9\le B_{\rm exact}$.
\end{proposition}
\begin{proof}
\textbf{Step 1: solve both radial boundary-value problems.}
For a radial field, $\Delta a=a''+2a'/s$. Substitution shows
$-\Delta[(1-s^2)/6]=1$ and $\Delta(1/s)=0$ for $s>0$.
Regularity at the origin and zero trace at $s=1$ give $u$.
On the annulus write $v=u+b(1/s-1)$; the condition $v(r)=0$ yields
$b=-r(1+r)/6$. Zero trace at one is automatic. Uniqueness of the weak
Dirichlet problem verifies these candidates are the desired solutions.

\textbf{Step 2: compute the exact difference energy.}
Inside the removed ball, $w=u$ and $|w'(s)|=s/3$, giving energy
$4\pi\int_0^r(s^2/9)s^2\,ds=4\pi r^5/45$.
Outside it, $w=r(1+r)(1/s-1)/6$. Its energy is
\[
 \frac{4\pi r^2(1+r)^2}{36}\int_r^1s^{-2}\,ds
 =\frac{\pi r(1-r)(1+r)^2}{9}.
\]
Adding the two energies proves the benefit formula by
\cref{eq:coolingprojection}. Direct radial integration gives
$J_0=4\pi\int_0^1(1-s^2)s^2/6\,ds=4\pi/45$.

\textbf{Step 3: independently check the lower bound.}
The minimum of $u$ on $\overline B_r$ is $(1-r^2)/6$.
Substituting into $4\pi rm^2$ gives the stated $B_{\rm floor}$.
Subtracting it from the exact expression leaves
\[
 B_{\rm exact}-B_{\rm floor}
 =\frac{\pi r^2(1-r)(1+r)^2}{9}+\frac{4\pi r^5}{45}\ge0.
\]
This checks both the capacity normalization and the sign of the guarantee
at finite radius, independently of the asymptotic theorem.
\end{proof}

\paragraph{Reproducible check.}
A numerical check evaluates the closed-form
state and compliance, checks the lower bounds over a fixed radius grid,
and runs the fixed-budget validator on this analytic geometry. It
records the seed, sample count, stopping policy, range bound,
confidence level, target, and actual exact benefit. It is a manufactured
check of the formulas and implementation, not a coverage proof or a
Bunny/housing performance measurement. The confidence statement comes from
the preceding derivation; numerical code remains subject to the explicitly
stated geometry, source-sampling, independence, and arithmetic assumptions.

\begin{table}[t]
\centering
\small
\caption{Manufactured minimum-benefit check on the unit ball, with a
preselected cavity radius $r=0.1$, containing radius $R=0.9$, and target
$10\%$ of the exact original compliance. The Monte Carlo rule allocates
$\delta=0.01$ to this one fixed test. These are numerical evaluations of
the proved formulas, not an empirical estimate of coverage.}
\label{tab:minimumbenefitcheck}
\begin{tabular}{lr}
\toprule
Quantity & Value\\
\midrule
Independent walks $N$ & $65{,}536$\\
Maximum steps $T$ / radius cap $s_*$ & $32$ / $0.35$\\
Exact center state $u(0)$ & $0.166667$\\
Truncated sample mean $\overline W$ & $0.162320$\\
One-sided sampling allowance & $0.003873$\\
State lower bound $L_u$ & $0.158447$\\
Source-only benefit bound / $J_0$ & $8.000\%$\\
Monte Carlo benefit bound / $J_0$ & $10.154\%$\\
Exact finite-hole benefit / $J_0$ & $13.614\%$\\
\bottomrule
\end{tabular}
\end{table}

For this fixed seed and budget, the source-only rule rejects the stipulated
target and the sampled rule accepts it (\cref{tab:minimumbenefitcheck}).
The exact energy formula agrees with independent radial quadrature to
$4.3\times10^{-16}$ in the seven tested cases. The deterministic inequalities
are also checked at $2{,}001$ radii below $R$; their validity for all radii
comes from the proofs, not from this finite grid. The sample budget was not
chosen as a solver-speed comparison, and one successful seed does not
establish practical tightness across application geometries.
